\documentclass[11pt,a4paper]{article}

\usepackage[utf8]{inputenc}
\usepackage[T1]{fontenc}
\usepackage[swedish]{babel}
\usepackage{amsmath, amssymb}
\usepackage{geometry}
\geometry{margin=2.5cm}
\usepackage{parskip}

\newcommand{\dd}{\mathrm{d}}

\title{Sammanfattning för poster}
\author{Projektarbete SF1693, KTH VT 2026}
\date{\today}

\begin{document}
\maketitle

\paragraph{OU-processen.}
Vi studerar Ornstein--Uhlenbeck-processen
$\dd X_t = \theta(\mu - X_t)\,\dd t + \sigma\,\dd W_t$ och härleder
dess analytiska lösning med integrerande faktor. Lösningen är
Gaussisk med medelvärde $\mu + (x_0-\mu)e^{-\theta t}$ och varians
$\sigma^2(1-e^{-2\theta t})/(2\theta)$, med stationär fördelning
$\mathcal{N}(\mu, \sigma^2/(2\theta))$.

\paragraph{Fokker--Planck.}
Tätheten $p(x,t)$ uppfyller PDE:n
$\partial_t p = -\partial_x[\theta(\mu-x)p] + \tfrac{1}{2}\sigma^2 \partial_{xx} p$.
Vi löser den numeriskt med finita differenser och jämför mot den
explicita Gaussiska referenstätheten.

\paragraph{Vasicek-modellen.}
Som finansiell tillämpning prissätter vi nollkupongsobligationer
under $\dd r_t = a(b-r_t)\,\dd t + \sigma\,\dd W_t$. Feynman--Kac
ger termstruktur-PDE:n, och en exponentiellt-affin ansats
$P = \exp(A - Br)$ reducerar PDE:n till ett ODE-system med sluten
lösning för $A(t,T)$ och $B(t,T)$.

\paragraph{Monte Carlo.}
Vi simulerar räntebanor med den exakta Gaussiska övergångsfördelningen,
trapetsintegrerar diskonteringen och medelvärdesbildar. Skattaren
stämmer med den slutna Vasicek-formeln inom statistisk osäkerhet vid
$N\sim 10^5$ banor.

\paragraph{Finita differenser för Vasicek.}
Termstruktur-PDE:n löses också direkt med Crank--Nicolson på ett
$(r,t)$-rutnät, med Dirichlet-rand från den slutna formeln. Med
parametrar skattade ur 10-åriga obligationsdata (Sverige och Indien)
återges priset $P(r,t)$ med maximalt fel $\sim 2\cdot 10^{-7}$.

\paragraph{Slutsats.}
OU-processen kopplar fysik och finans: en SDE med ren analytisk
lösning, en tillhörande PDE för tätheten, och en finansiell tillämpning
i Vasicek-modellen. Tre numeriska metoder --- FD för Fokker--Planck,
FD för obligationspriset och Monte Carlo --- bekräftar de slutna
formlerna.

\end{document}
